\appendix

\renewcommand{\thefigure}{\thesection.\arabic{figure}}
\renewcommand{\thetable}{\thesection.\arabic{table}}
\renewcommand{\theequation}{\thesection.\arabic{equation}}

\section{Correspondencia de Figuras y Cuadros}
\label{sec:anexos}

\subsection{Relación entre las Variables del Análisis}
\label{anx:fig_variables}

\begin{equation}
\text{Estabilidad}
+
\text{Eficiencia}
+
\text{Condiciones de Operación}
\longrightarrow
\text{Producción Útil}.
\end{equation}

\begin{center}
\small\hyperref[fig:variables]{Volver a la Figura \ref*{fig:variables}}
\end{center}


\subsection{Protocolo Utilizado para Comparar WT y ENG}
\label{anx:tab_protocolo}

\begin{table}[H]
\centering
\caption{Protocolo Utilizado para Comparar WT y ENG}
\caption*{\small\textit{Fuente:} Elaboración propia.}
\renewcommand{\arraystretch}{1.25}
\begin{tabularx}{0.96\textwidth}{|p{2.6cm}|p{4.1cm}|X|}
\hline
Prueba & Condición & Resultado Utilizado \\ \hline
Estabilidad & \(60\,^{\circ}\mathrm{C}\), 10 mM de sustrato &
\(t_{50}\) y actividad conservada en el tiempo. \\ \hline
Cinética & 0.5--80 mM de sustrato &
Eficiencia de conversión e inhibición por exceso de sustrato. \\ \hline
Respuesta Térmica & \(35--90\,^{\circ}\mathrm{C}\) &
Temperatura de máxima actividad. \\ \hline
Producción Base & \(60\,^{\circ}\mathrm{C}\), 40 mM, 8 horas &
Producto acumulado bajo una condición común. \\ \hline
Exploración Operativa & Misma malla de temperaturas y sustratos, 8 horas &
Mejor punto de producción encontrado para WT y ENG. \\ \hline
\end{tabularx}
\end{table}

\begin{center}
\small\hyperref[tab:protocolo_metodologia]{Volver al Cuadro \ref*{tab:protocolo_metodologia}}
\end{center}


\subsection{Lectura de los Índices}
\label{anx:tab_lectura_indices}

\begin{table}[H]
\centering
\caption{Lectura de los Índices}
\caption*{\small\textit{Fuente:} Elaboración propia.}
\renewcommand{\arraystretch}{1.25}
\begin{tabular}{|c|p{9.5cm}|}
\hline
Valor & Interpretación \\ \hline
\(>1\) & ENG mejora el resultado frente a WT. \\ \hline
\(=1\) & No existe diferencia entre ambas variantes. \\ \hline
\(<1\) & ENG obtiene un resultado inferior a WT. \\ \hline
\end{tabular}
\end{table}

\begin{center}
\small\hyperref[tab:lectura_indices]{Volver al Cuadro \ref*{tab:lectura_indices}}
\end{center}


\subsection{Flujo Metodológico Orientado a Decisión}
\label{anx:fig_flujo}

\begin{equation}
\text{Comparar WT--ENG}
\longrightarrow
\text{Medir}
\longrightarrow
\text{Identificar el Límite}
\longrightarrow
\text{Elegir el Punto Operativo}.
\end{equation}

\begin{center}
\small\hyperref[fig:flujo_metodologia]{Volver a la Figura \ref*{fig:flujo_metodologia}}
\end{center}


\subsection{Producción Acumulada Bajo la Condición Base}
\label{anx:fig_produccion}

\begin{figure}[H]
\centering
\begin{tabular}{|c|c|}
\hline
\parbox{0.46\textwidth}{\centering\small Caso 1: Procesamiento de Alimentos\\[0.12cm]\includegraphics[width=0.43\textwidth]{13.png}}
&
\parbox{0.46\textwidth}{\centering\small Caso 2: Bagazo de Caña\\[0.12cm]\includegraphics[width=0.43\textwidth]{14.png}}
\\ \hline
\parbox{0.46\textwidth}{\centering\small Caso 3: Procesamiento de Biomasa\\[0.12cm]\includegraphics[width=0.43\textwidth]{15.png}}
&
\parbox{0.46\textwidth}{\centering\small Caso 4: Síntesis de Glucósidos\\[0.12cm]\includegraphics[width=0.43\textwidth]{16.png}}
\\ \hline
\end{tabular}
\caption{Producción Acumulada Bajo la Condición Base}
\caption*{\small\textit{Fuente:} Elaboración propia.}
\end{figure}

\begin{center}
\small\hyperref[fig:produccion_casos]{Volver a la Figura \ref*{fig:produccion_casos}}
\end{center}


\subsection{Resumen Cuantitativo Bajo la Condición Base}
\label{anx:tab_resumen}

\begin{table}[H]
\centering
\caption{Resumen Cuantitativo Bajo la Condición Base}
\caption*{\small\textit{Fuente:} Elaboración propia.}
\renewcommand{\arraystretch}{1.25}
\setlength{\tabcolsep}{3.5pt}
\begin{tabularx}{\textwidth}{
|p{2.5cm}|
>{\centering\arraybackslash}p{2.0cm}|
>{\centering\arraybackslash}p{1.45cm}|
>{\centering\arraybackslash}p{1.45cm}|
>{\centering\arraybackslash}p{1.45cm}|
>{\centering\arraybackslash}p{2.0cm}|
X|}
\hline
Aplicación & \(t_{50}\) WT \(\rightarrow\) ENG & \(I_S\) & \(I_C\) & \(I_P\) &
\(T_{\mathrm{opt}}\) WT \(\rightarrow\) ENG & Lectura \\ \hline
Alimentos & 1.11 \(\rightarrow\) 3.53 h & 3.37 & 1.54 & 4.58 &
60 \(\rightarrow\) 65 \(^{\circ}\)C & La vida operativa impulsa el resultado. \\ \hline
Bagazo & 0.75 \(\rightarrow\) 7.02 h & 6.92 & 1.48 & 8.96 &
55 \(\rightarrow\) 65 \(^{\circ}\)C & Estabilidad y tolerancia al producto. \\ \hline
Biomasa & 0.99 \(\rightarrow\) 6.13 h & 4.86 & 1.58 & 5.41 &
65 \(\rightarrow\) 75 \(^{\circ}\)C & Mayor ventana térmica útil. \\ \hline
Glucósidos & 1.49 \(\rightarrow\) 5.43 h & 3.19 & 2.56 & 15.29 &
55 \(\rightarrow\) 60 \(^{\circ}\)C & Eficiencia y selectividad amplifican la mejora. \\ \hline
\end{tabularx}
\end{table}

\begin{center}
\small\hyperref[tab:resumen_resultados]{Volver al Cuadro \ref*{tab:resumen_resultados}}
\end{center}


\subsection{Síntesis Comparativa de los Cuatro Casos}
\label{anx:fig_sintesis}

\begin{figure}[H]
\centering
\begin{tabular}{|c|c|}
\hline
\parbox{0.46\textwidth}{\centering\small Impacto sobre Producción\\[0.12cm]\includegraphics[width=0.43\textwidth]{20.png}}
&
\parbox{0.46\textwidth}{\centering\small Matriz Estabilidad--Eficiencia\\[0.12cm]\includegraphics[width=0.43\textwidth]{21.png}}
\\ \hline
\end{tabular}
\caption{Síntesis Comparativa de los Cuatro Casos}
\caption*{\small\textit{Fuente:} Elaboración propia.}
\end{figure}

\begin{center}
\small\hyperref[fig:resumen_experimento]{Volver a la Figura \ref*{fig:resumen_experimento}}
\end{center}


\subsection{Restricciones No Lineales Observadas en los Cuatro Casos}
\label{anx:fig_restricciones}
\label{anx:fenomenos}

\begin{figure}[H]
\centering
\begin{tabular}{|c|c|}
\hline
\parbox{0.46\textwidth}{\centering\small Proteasa: Riesgo de Inactivación\\[0.12cm]\includegraphics[width=0.43\textwidth]{22.png}}
&
\parbox{0.46\textwidth}{\centering\small Bagazo: Inhibición por Producto\\[0.12cm]\includegraphics[width=0.43\textwidth]{23.png}}
\\ \hline
\parbox{0.46\textwidth}{\centering\small Biomasa: Activación Térmica Transitoria\\[0.12cm]\includegraphics[width=0.43\textwidth]{24.png}}
&
\parbox{0.46\textwidth}{\centering\small Glucósidos: Inhibición por Sustrato\\[0.12cm]\includegraphics[width=0.43\textwidth]{25.png}}
\\ \hline
\end{tabular}
\caption{Restricciones No Lineales Observadas en los Cuatro Casos}
\caption*{\small\textit{Fuente:} Elaboración propia.}
\end{figure}

\begin{center}
\small\hyperref[fig:fenomenos]{Volver a la Figura \ref*{fig:fenomenos}}
\end{center}


\subsection{Diagnóstico Industrial de los Cuatro Casos}
\label{anx:tab_diagnostico}

\begin{table}[H]
\centering
\caption{Diagnóstico Industrial de los Cuatro Casos}
\caption*{\small\textit{Fuente:} Elaboración propia.}
\renewcommand{\arraystretch}{1.25}
\begin{tabularx}{\textwidth}{|p{2.4cm}|p{3.2cm}|p{3.2cm}|X|}
\hline
Aplicación & Cuello de Botella & Evidencia Principal & Acción Prioritaria \\ \hline
Alimentos & Inactivación temprana & \(t_{50}\): 1.11 \(\rightarrow\) 3.53 h &
Explotar mayor vida operativa. \\ \hline
Bagazo & Inactivación e inhibición por glucosa &
\(t_{50}\): 0.75 \(\rightarrow\) 7.02 h & Priorizar estabilidad y tolerancia. \\ \hline
Biomasa & Ventana térmica &
\(T_{\mathrm{opt}}\): 65 \(\rightarrow\) 75 \(^{\circ}\)C &
Aprovechar temperaturas mayores. \\ \hline
Glucósidos & Inhibición y selectividad &
\(I_C=2.56\), \(I_P=15.29\) & Optimizar producto útil. \\ \hline
\end{tabularx}
\end{table}

\begin{center}
\small\hyperref[tab:diagnostico_industrial]{Volver al Cuadro \ref*{tab:diagnostico_industrial}}
\end{center}


\subsection{Traducción de la Mejora Molecular a Producción}
\label{anx:tab_traduccion}

\begin{table}[H]
\centering
\caption{Traducción de la Mejora Molecular a Producción}
\caption*{\small\textit{Fuente:} Elaboración propia.}
\renewcommand{\arraystretch}{1.25}
\begin{tabular}{|l|c|l|}
\hline
Aplicación & \(F_P\) & Lectura \\ \hline
Alimentos & 0.88 & Ligera atenuación por restricciones del proceso. \\ \hline
Bagazo de Caña & 0.87 & La inhibición absorbe parte de la mejora. \\ \hline
Biomasa & 0.70 & La estabilidad deja de ser el único límite. \\ \hline
Glucósidos & 1.87 & Selectividad y tolerancia amplifican el resultado. \\ \hline
\end{tabular}
\end{table}

\begin{center}
\small\hyperref[tab:factor_traduccion]{Volver al Cuadro \ref*{tab:factor_traduccion}}
\end{center}


\subsection{Puntos de Operación con Mayor Producto Dentro del Rango Evaluado}
\label{anx:tab_puntos}

\begin{table}[H]
\centering
\caption{Puntos de Operación con Mayor Producto Dentro del Rango Evaluado}
\caption*{\small\textit{Fuente:} Elaboración propia.}
\renewcommand{\arraystretch}{1.25}
\setlength{\tabcolsep}{3.5pt}
\begin{tabularx}{\textwidth}{
|p{2.1cm}|
>{\centering\arraybackslash}p{2.0cm}|
>{\centering\arraybackslash}p{1.5cm}|
>{\centering\arraybackslash}p{2.0cm}|
>{\centering\arraybackslash}p{1.5cm}|
>{\centering\arraybackslash}p{1.5cm}|
X|}
\hline
Aplicación & ENG: \(T\), Sustrato & Producto ENG &
WT: \(T\), Sustrato & Producto WT & Ventaja & Decisión \\ \hline
Alimentos & 60 \(^{\circ}\)C, 40 mM & 16.47 &
45 \(^{\circ}\)C, 40 mM & 5.93 & 2.78\(\times\) &
Usar ENG para extender producción. \\ \hline
Bagazo & 60 \(^{\circ}\)C, 40 mM & 12.70 &
45 \(^{\circ}\)C, 20 mM & 2.93 & 4.33\(\times\) &
Explotar estabilidad y tolerancia a glucosa. \\ \hline
Biomasa & 70 \(^{\circ}\)C, 80 mM & 36.04 &
50 \(^{\circ}\)C, 40 mM & 7.38 & 4.88\(\times\) &
Operar ENG a mayor temperatura y carga. \\ \hline
Glucósidos & 60 \(^{\circ}\)C, 40 mM & 9.49 &
45 \(^{\circ}\)C, 20 mM & 1.59 & 5.95\(\times\) &
Priorizar producto útil y selectividad. \\ \hline
\end{tabularx}
\end{table}

\begin{center}
\small\hyperref[tab:puntos_operacion]{Volver al Cuadro \ref*{tab:puntos_operacion}}
\end{center}


\section{Evidencia Gráfica Complementaria}
\label{sec:evidencia_complementaria}

\subsection{Estabilidad Térmica}
\label{anx:estabilidad}

\begin{figure}[H]
\centering
\begin{tabular}{|c|c|}
\hline
\parbox{0.46\textwidth}{\centering\small Procesamiento de Alimentos\\[0.12cm]\includegraphics[width=0.43\textwidth]{1.png}}
&
\parbox{0.46\textwidth}{\centering\small Bagazo de Caña\\[0.12cm]\includegraphics[width=0.43\textwidth]{2.png}}
\\ \hline
\parbox{0.46\textwidth}{\centering\small Procesamiento de Biomasa\\[0.12cm]\includegraphics[width=0.43\textwidth]{3.png}}
&
\parbox{0.46\textwidth}{\centering\small Síntesis de Glucósidos\\[0.12cm]\includegraphics[width=0.43\textwidth]{4.png}}
\\ \hline
\end{tabular}
\caption{Estabilidad Térmica}
\caption*{\small\textit{Fuente:} Elaboración propia.}
\end{figure}


\subsection{Cinética Frente a Concentración de Sustrato}
\label{anx:cinetica}

\begin{figure}[H]
\centering
\begin{tabular}{|c|c|}
\hline
\parbox{0.46\textwidth}{\centering\small Procesamiento de Alimentos\\[0.12cm]\includegraphics[width=0.43\textwidth]{5.png}}
&
\parbox{0.46\textwidth}{\centering\small Bagazo de Caña\\[0.12cm]\includegraphics[width=0.43\textwidth]{6.png}}
\\ \hline
\parbox{0.46\textwidth}{\centering\small Procesamiento de Biomasa\\[0.12cm]\includegraphics[width=0.43\textwidth]{7.png}}
&
\parbox{0.46\textwidth}{\centering\small Síntesis de Glucósidos\\[0.12cm]\includegraphics[width=0.43\textwidth]{8.png}}
\\ \hline
\end{tabular}
\caption{Cinética Frente a Concentración de Sustrato}
\caption*{\small\textit{Fuente:} Elaboración propia.}
\end{figure}


\subsection{Respuesta Térmica}
\label{anx:temperatura}

\begin{figure}[H]
\centering
\begin{tabular}{|c|c|}
\hline
\parbox{0.46\textwidth}{\centering\small Procesamiento de Alimentos\\[0.12cm]\includegraphics[width=0.43\textwidth]{9.png}}
&
\parbox{0.46\textwidth}{\centering\small Bagazo de Caña\\[0.12cm]\includegraphics[width=0.43\textwidth]{10.png}}
\\ \hline
\parbox{0.46\textwidth}{\centering\small Procesamiento de Biomasa\\[0.12cm]\includegraphics[width=0.43\textwidth]{11.png}}
&
\parbox{0.46\textwidth}{\centering\small Síntesis de Glucósidos\\[0.12cm]\includegraphics[width=0.43\textwidth]{12.png}}
\\ \hline
\end{tabular}
\caption{Respuesta Térmica}
\caption*{\small\textit{Fuente:} Elaboración propia.}
\end{figure}


\subsection{Selectividad en Síntesis de Glucósidos}
\label{anx:selectividad}

\begin{figure}[H]
\centering
\fbox{\includegraphics[width=0.65\textwidth]{26.png}}
\caption{Selectividad en Síntesis de Glucósidos}
\caption*{\small\textit{Fuente:} Elaboración propia.}
\end{figure}

\section{Reproducibilidad y Trazabilidad}
\label{sec:reproducibilidad}

\subsection{Condiciones Comunes}
\label{anx:condiciones_repro}
\label{subsec:condiciones_repro}

\begin{table}[H]
\centering
\caption{Condiciones Generales del Análisis}
\caption*{\small\textit{Fuente:} Elaboración propia.}
\label{tab:reproducibilidad}
\renewcommand{\arraystretch}{1.25}
\begin{tabular}{|p{6.0cm}|p{6.2cm}|}
\hline
Condición & Valor \\ \hline
Réplicas & 4 por condición \\ \hline
Evaluación de Estabilidad & \(60\,^{\circ}\mathrm{C}\) \\ \hline
Sustrato en Estabilidad & \(10\,\mathrm{mM}\) \\ \hline
Condición Base de Producción & \(60\,^{\circ}\mathrm{C}\), 40 mM, 8 horas \\ \hline
Rango de Sustrato & \(0.5--80\,\mathrm{mM}\) \\ \hline
Rango de Temperatura & \(35--90\,^{\circ}\mathrm{C}\), pasos de \(5\,^{\circ}\mathrm{C}\) \\ \hline
\end{tabular}
\end{table}

\subsection{Variables Registradas}
\label{anx:variables_repro}
\label{subsec:variables_repro}

\begin{enumerate}
    \item actividad residual durante exposición térmica;
    \item velocidad inicial frente a concentración de sustrato;
    \item actividad relativa frente a temperatura;
    \item producto acumulado durante ocho horas.
\end{enumerate}


\subsection{Indicadores de Comparación}
\label{anx:indicadores_repro}

\begin{equation}
I_S=
\frac{\text{Estabilidad ENG}}
{\text{Estabilidad WT}},
\qquad
I_C=
\frac{\text{Eficiencia ENG}}
{\text{Eficiencia WT}}
\end{equation}

\begin{equation}
I_P=
\frac{\text{Producto ENG}}
{\text{Producto WT}},
\qquad
F_P=
\frac{I_P}{I_S I_C}.
\end{equation}

\(I_S\), \(I_C\) e \(I_P\) mayores que 1 indican mejora frente a WT.
\(F_P<1\) indica que otras restricciones absorben parte de la mejora molecular;
\(F_P>1\) indica que otros factores amplifican el resultado.


\subsection{Funciones Utilizadas}
\label{anx:funciones}
\label{subsec:funciones_repro}

\begin{equation}
A(t)=
A_0
\exp\left[
-\left(
\frac{t}{\tau}
\right)^{\beta}
\right].
\end{equation}

\begin{equation}
v=
\frac{V_{\max}[S]}
{K_m+[S]+\dfrac{[S]^2}{K_i}}.
\end{equation}

\subsection{Trazabilidad del Análisis}
\label{anx:trazabilidad}
\label{subsec:trazabilidad}

\begin{equation}
\text{Medición}
\longrightarrow
\text{Comparación WT--ENG}
\longrightarrow
\text{Resultado Productivo}
\longrightarrow
\text{Punto de Operación}
\longrightarrow
\text{Decisión}.
\end{equation}
